\chapter[Referencial Teórico]{Referencial Teórico}
\label{cap:referencial}

Neste capítulo serão abordados os principais conceitos e assuntos associados aos estudos e as realizações importantes à melhor compreensão de seu contexto e o desenvolvimento condizente as demandas identificadas.

\section[Tutoria]{Tutoria}
\label{sec:tutoria}

A prática de uma pessoa mais experiente ensinar ou orientar uma pessoa mais novata no estudo ou no exercício de uma determinada atividade pode receber diversos nomes a depender do contexto: mentoria, tutoria ou monitoria.

Sendo assim, é importante explicar cada um dos conceitos citados para melhor discernimento do que o projeto visa alcançar, contribuindo ainda com o melhor esclarecimento sobre cada um deles e suas diferenças.

Mentor é um indivíduo experiente que tem por objetivo passar adiante suas experiências e vivências em uma determinada área de atuação. Normalmente age como uma referência, uma pessoa disposta a demonstrar tudo o que sabe, orientando, apontando direções e estabelecendo desafios que permitam ao aprendiz que está sendo orientado, denominado mentorado, desenvolver seu conhecimento a partir de suas experiências \cite{oliveira2020mentoria}.

A mentoria é um acompanhamento de longo prazo e marcada por uma relação de respeito e confiança entre o mentor e o mentorado, envolvendo aspectos de carreira e apoio psicológico \cite{oliveira2020mentoria}. Entretanto, o mentor tem um papel que ultrapassa a orientação para estudo, buscando também o desenvolvimento pessoal e psicossocial do mentorado \cite{botti2008preceptor}. Exemplos incluem a aquisição de competências socioemocionais e cognitivas e a melhoria da autoestima e autorregulação \cite{opp2024mentoria}.

Essa natureza pessoal da mentoria é presente desde o poema que deu origem a essa palavra: \textit{Odisseia} de Homero, em que Odisseu confia a Mentor, seu fiel e sábio amigo, a educação de seu filho Telêmaco antes de partir para a guerra de Tróia. Atena, deusa da sabedoria e da razão, assume a forma de Mentor e orienta Telêmaco no seu desenvolvimento físico, intelectual e espiritual \cite{semiao2009tutoria}.

Tutoria pode ser definida como uma relação de ajuda entre duas pessoas com diferentes níveis de conhecimento, competências e/ou experiências em que o tutorando, o indivíduo de menor capacidade, recebe apoio do tutor para superar dificuldades escolares ou acadêmicas \cite{opp2024mentoria}. O tutor não necessariamente não precisa ser um docente formal, podendo ser um familiar, estudante ou um voluntário \cite{topping2000tutoring}.

A tutoria existe desde a Antiguidade, onde a cultura de povos e valores como honra, justiça e honestidade eram passados pelos tutores para seus aprendizes. A partir das universidades na Idade Média a tutoria passou também a ser acadêmica, com tutores sendo alunos mais velhos que ensinavam habilidades como leitura e interpretação de texto para os estudantes mais jovens, sob supervisão de um docente \cite{geib2007tutoria}.

Comparado com o aprendizado tradicional com um professor, a tutoria pode fornecer mais experiência prática, mais atividades e variedade e também mais ajuda individualizada \cite{topping2000tutoring}. Embora parecidas, a tutoria e a mentoria se diferem em seus objetivos, visto que a tutoria visa ajudar na superação de dificuldades através do ensino de conteúdos ou competências específicas, sem buscar ativamente o desenvolvimento pessoal \cite{opp2024mentoria}. Entretanto, este ainda pode vir a ocorrer como consequência das interações, como aumento da autoconfiança e desenvolvimento de pensamento crítico, tanto para o tutor quanto para o tutorando \cite{lin2025peer}.

Monitoria é uma atividade que surgiu na Idade Média, e consiste em um estudante mais experiente de uma turma auxiliar o professor no ensino de colegas de classe mais novatos. Na metade do século XIV, os professores tinham quase sempre um monitor que os auxiliavam na escolarização. No século XVIII existia o Método Monitorial de Lancaster na Inglaterra, em que adolescentes, instruídos diretamente pelos mestres, atuavam como monitores ensinando outros adolescentes \cite{frison2016monitoria}. Nota-se que a monitoria continua a ser usada até os dias atuais.

No Brasil, o programa de monitoria está regulamentado pelas instituições de ensino superior e tem como finalidade promover a participação de estudantes com desempenho destacado em disciplinas, atuando como facilitadores da aprendizagem dos colegas e colaborando com o docente na construção de metodologias de ensino \cite{monitoria2025}.

Assim como na tutoria, ambos monitor e estudante se beneficiam da monitoria. O estudante recebe apoio no aprendizado do conteúdo da disciplina e o monitor desenvolve competências acadêmicas, sociais e profissionais, como comunicação, liderança, empatia e didática, além de seu engajamento em programas de monitoria poder abrir portas para novas oportunidades acadêmcias, como a participação em projetos de pesquisa e extensão \cite{monitoria2025}.

Diante do exposto, o modelo de interação pretendido pelo projeto é a tutoria, focada no aprendizado de conteúdos e habilidades específicas, não sendo o objetivo o desenvolvimento pessoal em conjunto com o aprendizado, inerente à mentoria, e não possuindo o caráter acadêmico intrínseco à monitoria. Entretanto, a tutoria pode evoluir para uma mentoria caso se forme um vínculo entre o tutor e o tutorado e o tutor queira o ajudar com seus problemas pessoais, passando a ser um mentor para o mesmo.

Neste trabalho, a aprendizagem é compreendida sob a ótica da Teoria da Aprendizagem Significativa de David Ausubel, que define aprendizagem como a organização e
integração de informações na estrutura cognitiva do indivíduo. De acordo com a teoria, a aprendizagem é significativa quando os novos conceitos são compreendidos e relacionados com conhecimentos prévios do indivíduo, criando conexões significativas e sendo mais provável de ser transferida para a memória de longo prazo, enquanto que uma aprendizagem mecânica ocorre quando o indivíduo não realiza tais conexões, simplesmente memorizando fatos e informações sem entender os princípios ou conexões subjacentes, podendo ocorrer descarte como conhecimento sem relevância, sem significado e sem sentido \cite{junior2023olhar}.

Considerando que o fator isolado mais importante para a aprendizagem é o que o aprendiz já sabe, a tutoria se apresenta como um mecanismo ideal para esse processo, pois permite que o tutor identifique os conhecimentos prévios do aprendiz e facilite a absorção da nova informação, promovendo a retenção do conhecimento e evitando a aprendizagem mecânica.

Adicionalmente, a dinâmica da tutoria é compreendida neste trabalho como um processo de aprendizagem colaborativa, fundamentado na perspectiva sociointeracionista de Lev Vygotsky. Segundo \citeonline{saul2025}, essa abordagem concebe o desenvolvimento cognitivo e o aprendizado como inerentemente sociais, mediados pela interação entre indivíduos com diferentes níveis de conhecimento. Nesse cenário, o aprendizado ocorre quando um indivíduo realiza uma tarefa ou atividade que está além de suas capacidades recebendo o suporte de uma pessoa ou recurso com mais conhecimento do que ele, denominado More Knowledgeable Other (MKO, ou "Outro Mais Conhecedor"). Através do compartilhamento do conhecimento, dicas e encorajamento, o aprendiz acessa sua Zona de Desenvolvimento Proximal (ZDP), justamente o espaço entre o que consegue fazer sozinho e com ajuda, internalizando novas habilidades e podendo se tornar um MKO para outros aprendizes.  

Para viabilizar a busca e o contato com os tutores de forma acessível, a solução proposta é o desenvolvimento de uma aplicação web. Esta plataforma servirá como um ambiente digital onde os usuários poderão não apenas buscar por especialistas, mas também realizar o contato inicial e agendamento das sessões de aprendizado.

\section[Aplicação Web]{Aplicação Web}
\label{sec:aw}

% O projeto visa desenvolver uma aplicação web.
Uma aplicação web é um programa computacional que é executado através de um navegador de Internet, também chamado de \textit{browser}, como Google Chrome, Mozilla Firefox e outros, com o objetivo de executar funções, como um formulário de contato ou uma calculadora. Exemplos de aplicações web são Google Docs e Microsoft Office 365 \cite{harwood2022internet}.

Segundo \citeonline{shklar2009web}, uma aplicação web é diferente de um \textit{website}. Um \textit{website} é essencialmente informativo, projetado para fornecer conteúdo de forma estática, semelhante a um folheto digital ou a um site de notícias, não necessariamente oferecendo funcionalidades. Essa distinção é reforçada por \citeonline{delaRosa2025webapp}, que destaca diferenças adicionais, como o nível de interatividade. \textit{Websites} tendem a oferecer menor interatividade, geralmente limitando-se a ações básicas, como navegação por \textit{links}, enquanto aplicações web permitem interações mais complexas, como envio de arquivos, preenchimento de formulários e edição de documentos, por serem projetados para oferecer funcionalidades para o usuário. 

Aplicações web são programas elaborados sob uma arquitetura cliente/servidor \cite{shklar2009web}. Segundo \citeonline{nyabuto2024architectural}, essa arquitetura funciona baseada no envio de requisições e respostas, em que o cliente inicia a comunicação enviando um pedido (requisição) e o servidor processa esse pedido e retorna uma resposta ao cliente. Na aplicação proposta, um exemplo de comunicação é o aprendiz buscar tutores, enviando uma requisição para o servidor, que a processa e retorna uma lista de tutores (resposta) para o aprendiz.

A interação do usuário com uma aplicação web pode ser considerada um diálogo. O cliente (usuário) fornece dados ou clica em elementos da interface, como um botão, para realizar uma ação. A aplicação então fornece uma resposta para a ação do usuário, como exibir dados em resposta a submissão de um formulário de pesquisa, que pode então gerar uma nova ação do usuário \cite{harwood2022internet}.

Para enriquecer a experiência na aplicação, serão implementadas duas estratégias principais. Primeiramente, um sistema de recomendação será desenvolvido para garantir que os usuários encontrem tutores compatíveis com seus interesses e necessidades. Adicionalmente, para promover o uso contínuo e manter os usuários engajados, serão aplicados elementos de Gamificação.

\subsection[Sistema de Recomendação]{Sistema de Recomendação}
\label{sec:recomendação}

Os Sistemas de Recomendação (SR) são algoritmos feitos para sugerir itens relevantes para usuários, sejam estes itens produtos, serviços, informações, etc. Normalmente os SR estão em aplicações web, aprendendo com interações anteriores dos usuários e/ou com atributos dos itens \cite{Aggarwal2016}. Esses sistemas ajudam os usuários a encontrarem itens de interesse em grandes conjuntos de dados, priorizando ou filtrando alternativas com base no histórico do usuário, preferências ou comportamentos similares de outros usuários \cite{Francesco2015}.

O Sistema de Recomendação Colaborativa (SRC) visa adquirir informações sobre o comportamento prévio ou as opiniões de um grupo de usuários existente para premeditar quais itens o usuário atual do sistema provavelmente irá gostar ou se interessar. Atualmente, esses tipos de sistemas são amplamente utilizados na indústria, especialmente como uma ferramenta em \textit{sites} de varejo \textit{online} para personalizar o conteúdo de acordo com as necessidades de um determinado cliente. Estas abordagens colaborativas usam uma matriz de classificações de itens de usuários como única entrada e normalmente produzem os seguintes tipos de saída: uma previsão (numérica) indicando até que ponto o usuário atual gostará ou não gostará de um determinado item e uma lista de \textbf{\textit{n}} itens recomendados. Essa lista dos \textbf{\textit{n}} melhores itens não deve, obviamente, conter itens que o usuário atual já tenha adquirido \cite{Jannach2011}. 

Nos Sistemas de Recomendação Baseados em Conteúdo (SRBC) os atributos descritivos dos itens são usados para fazer recomendações. O termo ``conteúdo'' refere-se a essas descrições. Nos métodos baseados em conteúdo, as classificações e o comportamento de aquisição dos usuários são combinados com as informações de conteúdo disponíveis nos itens. Nos métodos baseados em conteúdo as descrições de itens, que são rotuladas com classificações, são usadas como dados de treinamento para criar uma classificação específica do usuário ou um problema de modelagem de regressão. Para cada usuário, os documentos de treinamento correspondem às descrições dos itens que ele comprou ou classificou. A variável de classe (ou dependente) corresponde às classificações especificadas ou ao comportamento de aquisição. Esses documentos de treinamento são usados para criar um modelo de classificação ou regressão, que é específico para o usuário em questão (ou usuário ativo). Esse modelo específico do usuário é usado para prever se o indivíduo correspondente gostará de um item cuja classificação ou comportamento de aquisição é desconhecido \cite{Aggarwal2016}.

Os Sistemas de Recomendação Híbridos (SRH) surgem como uma opção interessante ao combinar diversas técnicas ou algoritmos, visando unir os pontos positivos de diferentes métodos, a fim de atender suas fraquezas e melhorar a eficiência das recomendações. A Recomendação Colaborativa é baseada na avaliação das preferências dos usuários e no comportamento coletivo da comunidade, a fim de produzir sugestões com base nas semelhanças entre perfis de usuários. Em contrapartida, os Sistemas de Recomendação Baseados em Conteúdo se valem de atributos próprios dos itens, como suas características estruturais e descrições em texto, para determinar a relevância pessoal de interações. Embora ambas as metodologias tenham benefícios reconhecidos, cada uma tem suas limitações específicas — como a limitação de dados na filtragem colaborativa quando pouco ou nenhum perfil de usuário para comparação e a insuficiência de variedade nas recomendações que se baseiam apenas no conteúdo \cite{Jannach2011}.

% Com base no exposto, e visando a maior compatibilidade possível entre usuários e tutores na aplicação proposta, optou-se pela implementação de um SR Híbrido. A escolha se justifica por combinar as forças das abordagens de conteúdo e colaborativa, mitigando suas fraquezas individuais.

% Em suma o componente baseado em conteúdo (SRBC) será responsável por recomendar tutores com base em atributos explícitos, como a área de conhecimento ou habilidade que o usuário busca (por exemplo, "Aulas de Violão"). Em paralelo, o componente de filtragem colaborativa (SRC) enriquecerá as sugestões ao analisar o comportamento de usuários com perfis semelhantes, recomendando tutores que foram bem avaliados ou procurados por eles. A união desses dois sistemas resulta em recomendações mais precisas e diversificadas.

\subsection[Gamificação]{Gamificação}
\label{sec:gamificacao}

Gamificação é a utilização de elementos de \textit{games} em atividades fora do contexto dos \textit{games}, com a finalidade de motivar os indivíduos à ação, auxiliar na solução de problemas e promover aprendizagens com o mesmo grau de envolvimento e motivação que normalmente são encontradas nos jogadores \cite{fardo2013gamificacao}.

A gamificação se apresenta como um fenômeno emergente e com muitas potencialidades de aplicação em diversos campos da atividade humana, pois a linguagem e metodologia dos \textit{games} são bastante populares e eficazes na resolução de problemas e aceitas naturalmente pelas atuais gerações, que cresceram interagindo com esse tipo de entretenimento. Assim, ela se justifica a partir de uma perspectiva sociocultural \cite{fardo2013gamificacao}.

Para compreender como os elementos de jogos influenciam o comportamento do usuário, \citeonline{chou2015actionable} propõe o framework \textit{Octalysis}. Diferente de abordagens que focam apenas na mecânica superficial, esta estrutura analisa a gamificação sob a ótica da motivação humana, identificando oito \textit{Core Drives} (Impulsos Centrais) que orientam nossas ações e decisões:

\begin{itemize}
    \item \textbf{Significado Épico e Chamado (Epic Meaning \& Calling):} Ocorre quando o usuário acredita estar engajado em algo maior que ele mesmo ou que foi ``escolhido'' para uma missão. É o pilar que motiva, por exemplo, contribuidores de projetos colaborativos como a Wikipédia, que dedicam tempo voluntário por uma causa maior.

    \item \textbf{Desenvolvimento e Realização (Development \& Accomplishment):} Trata-se do impulso interno de progredir, desenvolver habilidades e superar desafios. Este é o motor principal por trás da maioria dos sistemas de pontos, medalhas e rankings, pois os usuários sentem satisfação ao visualizar seu crescimento e completar tarefas difíceis.

    \item \textbf{Empoderamento da Criatividade e Feedback (Empowerment of Creativity \& Feedback):} Envolve engajar o usuário em um processo criativo onde ele precisa descobrir novas combinações e estratégias. O engajamento surge da autonomia para testar diferentes abordagens e receber feedback imediato sobre o resultado de suas escolhas.

    \item \textbf{Propriedade e Posse (Ownership \& Possession):} Baseia-se na motivação gerada quando o usuário sente que possui algo. Ao sentir-se dono de um item (virtual ou não), o usuário tende a querer melhorá-lo, protegê-lo e acumular mais. Este impulso justifica o uso de bens virtuais, coleções e personalização de perfis.

    \item \textbf{Influência Social e Relacionamento (Social Influence \& Relatedness):} Incorpora todos os elementos sociais que motivam as pessoas, como mentoria, aceitação social, competição e companheirismo. A observação do comportamento de outros usuários e o desejo de pertencimento são fortes catalisadores de ação.

    \item \textbf{Escassez e Impaciência (Scarcity \& Impatience):} É o impulso de desejar algo simplesmente porque é raro, exclusivo ou difícil de obter imediatamente. A limitação de tempo ou de acesso cria um senso de urgência que pode acelerar a tomada de decisão do usuário.

    \item \textbf{Imprevisibilidade e Curiosidade (Unpredictability \& Curiosity):} Relaciona-se ao engajamento gerado pela incerteza do que virá a seguir. A curiosidade cognitiva motiva o usuário a continuar interagindo com o sistema para descobrir novos resultados ou conteúdos.

    \item \textbf{Perda e Aversão (Loss \& Avoidance):} É a motivação baseada no medo de perder algo conquistado ou de que um evento negativo ocorra. No contexto de gamificação, pode ser utilizado para incentivar a continuidade do uso, evitando que o progresso ou status acumulado seja perdido.
\end{itemize}

A aplicação desses impulsionadores permite que sistemas gamificados, como a plataforma proposta neste trabalho, utilizem elementos como pontos e conquistas não apenas como recompensas arbitrárias, mas como ferramentas estruturadas para satisfazer necessidades psicológicas de competência e progresso.

Um sistema gamificado é composto de elementos que podem ser encontrados em \textit{games}. Alguns destes elementos podem ser definidos como Pontos, Níveis, Rankings, Desafios/Missões e Conquistas \cite{klock2014gamificacaoAVA}.

Os pontos são abertos, diretos e motivacionais, permitindo a utilização de vários tipos diferentes de pontuação, dependendo do objetivo. Alguns dos tipos de pontos podem ser descritos como pontos de experiência, pontos de habilidade e pontos de reputação \cite{zichermann2011gamificacao}.

Os níveis podem indicar o progresso do usuário dentro do sistema gamificado. Os níveis podem ser definidos como níveis de jogo, níveis de dificuldade e níveis do usuário \cite{zichermann2011gamificacao}.

Os Rankings são a comparação entre usuários envolvidos, a fim de visualizar suas progressões dentro do ambiente e inserir um senso de competição entre eles \cite{klock2014gamificacaoAVA}. Os desafios e missões norteiam os usuários sobre as atividades que devem ser feitas dentro do sistema \cite{fadel2014gamificacao}.

Conquistas são uma outra forma de pontos, possuindo uma representação visual de alguma realização do usuário no sistema \cite{hunter2012gamificacao}.

% Diante dos conceitos expostos, a estratégia de gamificação da aplicação será implementada com o objetivo de motivar e engajar os usuários. Para isso, serão utilizados os elementos de pontos, níveis, conquistas e rankings, que se interligam da seguinte forma: os usuários recebem conquistas ao realizarem ações relevantes na plataforma, como avaliar um tutor ou completar uma sessão de tutoria. Essas conquistas recompensam o usuário com pontos, cuja quantidade varia conforme a dificuldade da tarefa. 

% O acúmulo de pontos, por sua vez, determina o nível do usuário, que é representado por títulos decorativos em seu perfil (como "Iniciante" ou "Grão-Mestre!"). O propósito desses títulos é oferecer um senso de progresso e status, incentivando a participação contínua. Por fim, para estimular uma competição saudável e dar visibilidade aos tutores mais dedicados, a aplicação contará com um ranking que destacará os melhores classificados com base em uma combinação de suas avaliações e do número de tutorias realizadas por área de conhecimento.

A proposta, que será apresentada mais detalhadamente no Capítulo 3, é desenvolver uma aplicação web (programa computacional executado através do navegador de Internet) para facilitar o contato entre um tutor, que busca ajudar alguém a aprender ou se aperfeiçoar em uma habilidade, e um aprendiz que busca ajuda. A aplicação contará com um sistema de recomendação para sugerir tutores relevantes a cada usuário e gamificação (elementos de jogos digitais, como pontos e níveis) para manter tanto tutores quanto aprendizes engajados com a plataforma.